\begin{frame}
	\fta{Introduction}

	\renewcommand{\figurename}{Figure}

	\s{        
		While the calculations for series connections without parallel-connected 
		components are relatively straightforward, the analysis of 
		direct current networks with nodes that have n branches, where the
		number of branches is $n \geq 2$, quickly becomes confusing as the number of nodes increases.
		A multi-layered consideration of the DC network is 
		necessary here. The following methods are available for calculating DC network:
		}

	\b{
		The “Advanced DC Networks” module covers the following topics:
		}

	\begin{itemize}
		\item Node and mesh rule
		\item Superposition
		\item Node potential analysis
		\item Mesh current analysis
	\end{itemize}

	\s{
		The analysis of electrical networks with more than one current 
        or voltage source is also not trivial. Determining all currents and 
        voltages in the electrical network shown in Figure \ref{BildEinleitungsnetzwerk} 
        is not possible by calculating series and parallel connections alone. 
		}
 
	\fu{
		\resizebox{.5\textwidth}{!}{%
		\input{Tikz/src/beispielnetzwerk.tex}%
		}
	}{{\bf Example Network.} Direct current electrical network with a plurality of resistors,
    current sources and voltage sources. \label{BildEinleitungsnetzwerk}}
\end{frame}